\section{Introduction}
\label{sec:intro}

Incompressible viscous flow is governed by the Navier--Stokes equations, a coupled, nonlinear system in
which the pressure acts as a Lagrange multiplier for the divergence-free constraint. Finite-volume,
finite-element and spectral solvers have been refined for this system over five decades; for a fixed
geometry and a fixed set of parameters they are fast, robust and come with a mature theory of discretisation
error. Physics-informed neural networks (PINNs)~\cite{raissi2019pinn} approach the same equations from a
different direction. A neural network represents the solution as a smooth function of space and time, the
residuals of the differential equations are evaluated at scattered points by automatic differentiation, and
the network parameters are chosen to make these residuals, together with boundary and initial mismatches,
small in a least-squares sense. No mesh is built, derivatives of the representation are exact, and measured
data can be added to the same objective, which makes the approach attractive for flow reconstruction from
sparse observations~\cite{raissi2020hfm} and for surrogates that are queried continuously.

The price of this flexibility is an optimisation problem that is often badly conditioned. Plain
multilayer perceptrons fit smooth components first and resolve sharp features late or never; the loss is a
sum of terms whose gradients differ by orders of magnitude; the solution of an unsteady problem can be
learnt out of temporal order; and convection-dominated flows make the landscape stiff enough that training
collapses onto trivial or steady states~\cite{wang2021gradient,wang2022ntk,krishnapriyan2021failure}. A
sizeable literature proposes remedies: random Fourier features against spectral
bias~\cite{tancik2020fourier}, gated (modified) multilayer perceptrons and gradient-based loss
balancing~\cite{wang2021gradient}, weights derived from the neural tangent kernel~\cite{wang2022ntk},
causal weighting in time~\cite{wang2024causality}, residual-based adaptive sampling~\cite{wu2023rad}, random
weight factorisation~\cite{wang2022rwf}, adaptive residual architectures~\cite{wang2024piratenets}, exactly
divergence-free output parametrisations~\cite{jin2021nsfnets}, boundary conditions built into the
ansatz~\cite{lagaris1998,sukumar2022distance}, separable networks that make three-dimensional unsteady
problems affordable~\cite{cho2023spinn}, domain decomposition~\cite{jagtap2020xpinn} and operator
learning over families of problems~\cite{lu2021deeponet,wang2021pideeponet,li2021fno}. Practical
recommendations that combine several of these are collected in~\cite{wang2023expert}.

Individual remedies are usually reported on the problem that motivated them, and the combined recipe is
seldom tested against independent references with the failures documented. This paper does that for one
code base and one set of choices. We implemented the components above in JAX~\cite{jax2018github}, following
the design of JAX-PI~\cite{jaxpi} and SPINN~\cite{spinncode} but with our own residual and memory handling,
and applied them to a ladder of standard problems: the two-dimensional Taylor--Green vortex with its exact
solution, the lid-driven cavity against Ghia et al.~\cite{ghia1982}, the DFG 2D-2 cylinder benchmark of
Sch\"afer and Turek~\cite{schaefer1996}, the three-dimensional Taylor--Green vortex at $\Rey=1600$, an inverse
problem on the cylinder-wake data of~\cite{raissi2019pinn}, and a physics-informed DeepONet over the cavity
Reynolds number. Each benchmark has an acceptance gate fixed before training.

The contributions are the following.
\begin{enumerate}
\item A single, tested implementation in which every residual is assembled from forward-mode directional
  derivatives and evaluated in rematerialised chunks, so that third-order derivatives of a six-layer network
  on $\sim\!10^4$ points per step fit in about 1.6~GB of GPU memory.
\item An independent finite-difference reference for the regularised-lid cavity. It shows that a
  grid-converged solution of the benchmark as posed differs from the Ghia $v$ table by 2.3\%, so a 2\% gate
  against that table cannot be met even by an exact solution; the PINN is instead graded against the converged
  field of the same problem, which it reproduces to 0.27\%.
\item An attribution of the cavity failures. Imposing the regularised lid exactly is incompatible with
  continuity at the two top corners, which leaves an irreducible divergence residual; gradient-norm loss
  balancing then drives the weight of a softly imposed lid up until the incompatibility returns. Both effects
  are shown in controlled runs.
\item Results for all six problems with the cost of each run (wall-clock, step time, memory, inference
  throughput), including the gates that were missed and the budgets that had to be reduced.
\end{enumerate}

Section~\ref{sec:method} describes the formulations, architectures, losses and training procedure;
Section~\ref{sec:setup} the benchmarks, reference data and protocol; Section~\ref{sec:results} the results;
Section~\ref{sec:discussion} discusses them against classical solvers and lists the limitations.
